\documentclass[otchet]{SCWorks}
% Тип обучения (одно из значений):
%    bachelor   - бакалавриат (по умолчанию)
%    spec       - специальность
%    master     - магистратура
% Форма обучения (одно из значений):
%    och        - очное (по умолчанию)
%    zaoch      - заочное
% Тип работы (одно из значений):
%    coursework - курсовая работа (по умолчанию)
%    referat    - реферат
%  * otchet     - универсальный отчет
%  * nirjournal - журнал НИР
%  * digital    - итоговая работа для цифровой кафедры
%    diploma    - дипломная работа
%    pract      - отчет о научно-исследовательской работе
%    autoref    - автореферат выпускной работы
%    assignment - задание на выпускную квалификационную работу
%    review     - отзыв руководителя
%    critique   - рецензия на выпускную работу

% * Добавлены вручную. За вопросами к @mchernigin
\usepackage{preamble}

\begin{document}

% Кафедра (в родительном падеже)
\chair{математической кибернетики и компьютерных наук}

% Тема работы
\title{Анализ сложности алгоритма Бойера-Мура}

% Курс
\course{2}

% Группа
\group{211}

% Факультет (в родительном падеже) (по умолчанию "факультета КНиИТ")
% \department{факультета КНиИТ}

% Специальность/направление код - наименование
 \napravlenie{02.03.02 "--- Фундаментальная информатика и информационные технологии}
% \napravlenie{02.03.01 "--- Математическое обеспечение и администрирование информационных систем}
% \napravlenie{09.03.01 "--- Информатика и вычислительная техника}
%\napravlenie{09.03.04 "--- Программная инженерия}
% \napravlenie{10.05.01 "--- Компьютерная безопасность}

% Для студентки. Для работы студента следующая команда не нужна.
% \studenttitle{студентки}

% Фамилия, имя, отчество в родительном падеже
\author{Бурова Арсения Александровича}

% Заведующий кафедрой 
\chtitle{доцент, к.\,ф.-м.\,н.}
\chname{С.\,В.\,Миронов}

% Руководитель ДПП ПП для цифровой кафедры (перекрывает заведующего кафедры)
% \chpretitle{
%     заведующий кафедрой математических основ информатики и олимпиадного\\
%     программирования на базе МАОУ <<Ф"=Т лицей №1>>
% }
% \chtitle{г. Саратов, к.\,ф.-м.\,н., доцент}
% \chname{Кондратова\, Ю.\,Н.}

% Научный руководитель (для реферата преподаватель проверяющий работу)
\satitle{доцент, к.\,ф.-м.\,н.} %должность, степень, звание
\saname{Ю.\,Н.\,Кондратова}

% Руководитель практики от организации (руководитель для цифровой кафедры)
\patitle{доцент, к.\,ф.-м.\,н.}
\paname{С.\,В.\,Миронов}

% Руководитель НИР
\nirtitle{доцент, к.\,п.\,н.} % степень, звание
\nirname{В.\,А.\,Векслер}

% Семестр (только для практики, для остальных типов работ не используется)
\term{2}

% Наименование практики (только для практики, для остальных типов работ не
% используется)
\practtype{учебная}

% Продолжительность практики (количество недель) (только для практики, для
% остальных типов работ не используется)
\duration{2}

% Даты начала и окончания практики (только для практики, для остальных типов
% работ не используется)
\practStart{01.07.2024}
\practFinish{13.01.2024}

% Год выполнения отчета
\date{2026}

\maketitle

% Включение нумерации рисунков, формул и таблиц по разделам (по умолчанию -
% нумерация сквозная) (допускается оба вида нумерации)
\secNumbering

\tableofcontents

% Раздел "Обозначения и сокращения". Может отсутствовать в работе
% \abbreviations
% \begin{description}
%     \item ... "--- ...
%     \item ... "--- ...
% \end{description}

% Раздел "Определения". Может отсутствовать в работе
% \definitions

% Раздел "Определения, обозначения и сокращения". Может отсутствовать в работе.
% Если присутствует, то заменяет собой разделы "Обозначения и сокращения" и
% "Определения"
% \defabbr

\section{Реализация алгоритма}
\begin{Verbatim}[
  numbers=left,
  breaklines=true,
  breakanywhere=true
]
vector<int> buildBadCharTable(const string& pattern) {
    const int ALPHABET_SIZE = 256;
    vector<int> badCharTable(ALPHABET_SIZE, -1);

    for (int i = 0; i < pattern.size(); i++) {
        unsigned char currentChar = pattern[i];  //явно преобразуем в unsigned char
        int asciiCode = currentChar;              //получаем ascii код от 0 до 255
        badCharTable[asciiCode] = i;              //используем ascii код как индекс
    }

    return badCharTable;
}
\end{Verbatim}

\begin{Verbatim}[
  numbers=left,
  breaklines=true,
  breakanywhere=true
]
vector<int> boyerMoore(const string& text, const string& pattern) {
    vector<int> result;  // вектор для хранения позиций совпадений
    int n = text.size();  // длина текста
    int m = pattern.size();  // длина паттерна

    if (m == 0 || n < m) return result;  // если паттерн пуст или длиннее текста

    vector<int> badChar = buildBadCharTable(pattern);  // строим таблицу плохих символов
    int shift = 0;  // текущее смещение паттерна относительно текста

    while (shift <= n - m) {  // пока паттерн помещается в оставшийся текст
        int j = m - 1;  // начинаем сравнение с конца паттерна

        // сравниваем символы справа налево
        while (j >= 0 && pattern[j] == text[shift + j]) {
            j--;  // при совпадении двигаемся к началу
        }

        // если прошли весь паттерн, значит нашли совпадение
        if (j < 0) {
            result.push_back(shift);  // сохраняем позицию совпадения

            // вычисляем сдвиг для поиска следующего вхождения
            if (shift + m < n) {
                // получаем символ из текста, следующий сразу после паттерна
                char nextCharAfterPattern = text[shift + m];
                // преобразуем в unsigned char для безопасного индекса
                unsigned char unsignedNextChar = (unsigned char)nextCharAfterPattern;
                // получаем последнюю позицию этого символа в паттерне
                int lastPosInPattern = badChar[unsignedNextChar];
                // сдвигаем так чтобы этот символ совпал с его последним вхождением
                shift += m - lastPosInPattern;
            }
            else {
                shift += 1;  // если паттерн в конце текста, сдвигаем на 1
            }
        }
        else {
            // вычисляем сдвиг по эвристике плохого символа
            // получаем символ из текста на позиции несовпадения
            char mismatchChar = text[shift + j];
            unsigned char unsignedMismatchChar = (unsigned char)mismatchChar;
            int lastPos = badChar[unsignedMismatchChar]; // получаем последнюю позицию этого символа в паттерне
            // вычисляем на сколько можно сдвинуться
            int badCharShift = j - lastPos;
            // сдвигаем минимум на 1, чтобы не зациклиться
            shift += max(1, badCharShift);
        }
    }

    return result;
}
\end{Verbatim}

\newpage

\section{Анализ алгоритма}

Функция \texttt{buildBadCharTable} получает строку \texttt{pattern} длины $m$.

Создание вектора \texttt{badCharTable} размера 256 и его инициализация значением $-1$ выполняется за время $O(1)$, так как размер алфавита фиксирован.

\begin{verbatim}
vector<int> badCharTable(ALPHABET_SIZE, -1);
\end{verbatim}

Основной цикл выполняется для $i$ от $0$ до $m-1$, то есть $O(m)$ итераций.

\begin{verbatim}
for (int i = 0; i < pattern.size(); i++) {
    unsigned char currentChar = pattern[i];
    int asciiCode = currentChar;
    badCharTable[asciiCode] = i;
}
\end{verbatim}

Внутри цикла все операции выполняются за $O(1)$.

Общая сложность построения таблицы плохих символов:

$$
T_{\text{table}} = O(1) + O(m) \cdot O(1) = O(m)
$$

Функция \texttt{boyerMoore} получает строку \texttt{text} длины $n$ и строку \texttt{pattern} длины $m$.

Инициализация переменных выполняется за $O(1)$.

Построение таблицы плохих символов — $O(m)$.

Основной цикл \texttt{while} выполняется, пока $shift \le n - m$.

Количество итераций основного цикла соответствует количеству сдвигов шаблона относительно текста. В худшем случае количество сдвигов может достигать $O(n)$.

Внутри основного цикла находится внутренний цикл \texttt{while}, который сравнивает символы шаблона и текста справа налево. В худшем случае при каждом сдвиге может потребоваться $O(m)$ сравнений.

Операции внутри условных операторов (добавление в вектор, вычисление сдвига по таблице плохих символов) выполняются за $O(1)$.

Пусть:
\begin{itemize}
\item $n$ — длина текста,
\item $m$ — длина шаблона.
\end{itemize}

\textbf{Худший случай:}

При каждом сдвиге может потребоваться $O(m)$ сравнений во внутреннем цикле. Количество сдвигов в худшем случае составляет $O(n)$. Таким образом:

$$
T_{\text{worst}} = O(m) + O(n) \cdot O(m) = O(n \cdot m)
$$

\textbf{Лучший случай:}

При каждом сдвиге эвристика плохого символа позволяет сдвинуть шаблон на значительное расстояние. Количество сдвигов ограничено $O\left(\frac{n}{m}\right)$. При каждом сдвиге выполняется $O(1)$ сравнений:

$$
T_{\text{best}} = O(m) + O\left(\frac{n}{m}\right) \cdot O(1) = O(m) + O\left(\frac{n}{m}\right)
$$

\end{document}
